\subsection*{6.2 Stokes' Theorem}

\subsubsection*{6.2.1 Conceptual Explanation}

\begin{conceptbox}
    Stokes' Theorem states that the circulation of a vector field along a closed curve $C$ is equal to the total curl (rotation) of the field over the surface $S$ bounded by $C$:
    \begin{equation*}
        \oint_C \mathbf{F} \cdot d\mathbf{r}
        =
        \iint_S (\nabla \times \mathbf{F}) \cdot d\mathbf{S}
    \end{equation*}
\end{conceptbox}

In simple terms:
\begin{itemize}
    \item The left side measures how much the vector field “flows around” the boundary $C$
    \item The right side measures how much the field “rotates” across the entire surface $S$
\end{itemize}

\textbf{Key Idea:}
\[
    \boxed{\text{Total circulation around a boundary = total curl inside the surface}}
\]

This is extremely useful because:
\begin{itemize}
    \item A complicated line integral can be replaced with an easier surface integral
    \item We only need the curl, not the full path traversal
\end{itemize}

\begin{examplebox}
    Evaluate $\oint_C \mathbf{F}\cdot d\mathbf{r}$ where $\mathbf{F}=\langle -y,x,z\rangle$ and $C$ is the boundary of the disk $x^2+y^2\le1$ in the plane $z=1$.

    \textbf{Step 1: Understand the geometry}

    - The curve $C$ is a circle of radius 1
    - It lies in the plane $z=1$
    - The surface $S$ is the flat disk bounded by this circle

    Instead of computing the line integral directly (which would be complicated), we apply Stokes’ Theorem and convert it into a surface integral.

    \textbf{Step 2: Compute the curl of the vector field}

    Given:
    \[
        \mathbf{F} = \langle -y,\; x,\; z \rangle
    \]

    Compute:
    \[
        \nabla \times \mathbf{F}
        =
        \begin{vmatrix}
            \mathbf{i} & \mathbf{j} & \mathbf{k} \\
            \partial_x & \partial_y & \partial_z \\
            -y         & x          & z
        \end{vmatrix}
    \]

    Expand determinant:

    \begin{align*}
        \nabla \times \mathbf{F}
         & = \mathbf{i}\left(\frac{\partial z}{\partial y} - \frac{\partial x}{\partial z}\right)
        - \mathbf{j}\left(\frac{\partial z}{\partial x} - \frac{\partial (-y)}{\partial z}\right)
        + \mathbf{k}\left(\frac{\partial x}{\partial x} - \frac{\partial (-y)}{\partial y}\right)
    \end{align*}

    Compute each term:

    \begin{align*}
        \frac{\partial z}{\partial y} & = 0, \quad \frac{\partial x}{\partial z} = 0     \\
        \frac{\partial z}{\partial x} & = 0, \quad \frac{\partial (-y)}{\partial z} = 0  \\
        \frac{\partial x}{\partial x} & = 1, \quad \frac{\partial (-y)}{\partial y} = -1
    \end{align*}

    So:
    \begin{align*}
        \nabla \times \mathbf{F}
         & = \langle 0,\; 0,\; 1 - (-1) \rangle \\
         & = \langle 0,\; 0,\; 2 \rangle
    \end{align*}

    \textbf{Step 3: Determine the surface normal}

    The surface is a flat disk in the plane $z=1$.

    - The upward normal vector is:
    \[
        \mathbf{n} = \langle 0,0,1 \rangle
    \]

    Thus:
    \[
        d\mathbf{S} = \mathbf{n}\, dA = \langle 0,0,1 \rangle dA
    \]

    \textbf{Step 4: Compute the dot product}

    \begin{align*}
        (\nabla \times \mathbf{F}) \cdot d\mathbf{S}
         & = \langle 0,0,2 \rangle \cdot \langle 0,0,1 \rangle dA \\
         & = 2\, dA
    \end{align*}

    So the surface integral becomes:
    \[
        \iint_S 2\, dA
    \]

    \textbf{Step 5: Evaluate the surface integral}

    Since the integrand is constant (2), we can factor it out:
    \begin{align*}
        \iint_S 2\, dA = 2 \iint_S dA
    \end{align*}

    But:
    \[
        \iint_S dA = \text{Area of the disk}
    \]

    The disk has radius 1:
    \[
        \text{Area} = \pi (1)^2 = \pi
    \]

    Thus:
    \begin{align*}
        \iint_S 2\, dA = 2\pi
    \end{align*}

    \textbf{Final Answer:}
    \[
        \boxed{2\pi}
    \]
\end{examplebox}